这里把“似然的计算”和“用什么目标训练”分开一下。
前面的不等式说明：在相应路径正则性与权重设定下，控制分数误差，
可以控制反向 SDE 生成分布的误差；它并没有直接给出每个输入样本的概率密度值。

若使用后面要介绍的概率流 ODE，记其速度为 $v_\theta(x,t)$。
一个小时间步的映射是 $x\mapsto x+h v_\theta(x,t)$，
Jacobian 为 $I+hD_xv_\theta$。利用
$\log\det(I+hD_xv_\theta)=h\operatorname{tr}(D_xv_\theta)+O(h^2)$，
变量替换公式沿轨迹给出
\[
\frac{\mathrm d}{\mathrm dt}\log p_t^\theta(x_t)
=-\nabla\cdot v_\theta(x_t,t).
\]
所以从一个观测 $x_0$ 正向积分到 $x_T$ 后，
\[
\log p_0^\theta(x_0)=\log\pi(x_T)
 +\int_0^T\nabla\cdot v_\theta(x_t,t)\,\mathrm dt.
\]
先验密度可以计算，积分项又由网络向量场决定，这就是计算模型似然的路径。
网络不是完美分数时，仍可在可逆性等条件下定义它所诱导的 ODE 模型密度；
但这个密度是否逼近真实数据，是另一个学习问题。
实际数值积分和随机迹估计也会有误差。

因此，一阶分数训练不是“不能计算 ODE 的似然”，
而是不能仅凭一个很小的一阶训练误差，就无条件保证 ODE 模型的似然误差也很小。
有关高阶分数误差的结果是在附加假设下分析这种控制关系，
不是要求每个可以计算似然的模型都必须另训二阶和三阶网络。\cite{sde,theorybook}
